\section{Limitations and scope}

The main interface supplies the current symbolic feasible-action menu. This isolates action selection but does not test unrestricted proposal or language generation. The results establish conditional planning among executable phrases. They do not establish an autonomous language agent.

The completed evidence comes from a stylized arithmetic environment. Templated language and exact counterfactual transitions are valuable for mechanism analysis, but they limit external validity. Faithful iGSM and at least one non-arithmetic domain are required before submission.

Several mechanism controls use one seed. Direct ranking currently matches full JEPA in that setting. Detached geometry, geometry-only planning, local metric correlations, and controlled semantic interventions need five-seed completion before they support a central claim.

The repaired deeper-planning protocol uses a future feasible-action tree beyond depth one. It is candidate privileged and still shows declining performance. A deployable test-time-scaling result requires a shared learned proposal interface or full-catalogue search, paired fixed-work evaluation, and calibrated FLOPs for recurrent baselines.

The present model is trained from scratch on a controlled language. Its findings do not directly establish how pretrained large language models would organize open-domain planning states. Conversely, the controlled setting enables exact causal comparisons that would be difficult to interpret in an open-ended agent benchmark.
